\begin{frame}[fragile]{세 열을 읽으면 RNN 학습의 ``운명''}
  \begin{itemize}
    \item \kb{첫째 열 (승수 0.45 $< 1$)}: 전형적 \textbf{소실}.
      $T=50$이면 $10^{-17}$으로, \texttt{float32}의 표현 하한
      ($\sim 10^{-38}$) 근처까지 밀림 --- ``100단어 전의 주어''의
      그래디언트는 \textbf{수치적으로 0과 구별 불가능}, 그 자리의
      가중치는 \textbf{영원히 갱신되지 못}함
    \item \kb{둘째 열 (승수 0.75)}: 같은 소실이지만 \textbf{16,000배
      느리다} --- $w_{hh}$를 0.9에서 1.5로 올리자 소실은
      \textbf{사라지지 않고 속도만 줄었다}. ``게이트 없이 소실을
      극복하려 하면 \textbf{소실$\leftrightarrow$폭발 사이 줄타기}에
      불과''하다는 뜻
    \item \kb{셋째 열 (승수 1.5 $> 1$)}: \textbf{폭발}. $T=20$에서
      이미 $2.2 \times 10^{3}$, $T=50$이 넘으면 $10^{10}$ 대로
      \texttt{NaN} --- 11.2절의 \verb|np.clip| 그래디언트
      클리핑이 바로 이 열을 막기 위한 \textbf{급조된 처방}
      (Pascanu et al., 2013)
  \end{itemize}
  \vskip0.4em
  \kq{승수 1 하나를 ``조금만'' 넘기면(1에 가까이) 소실이 느려지고,
  크게 넘기면 폭발이 된다. 어떤 초기화를 잡아도 \textbf{지수함수
  라는 형 자체}는 바뀌지 않는다 --- 이 형을 바꾸는 것이 곧
  \textbf{게이트의 임무}.}
\end{frame}
